\documentclass[10pt,a4paper]{amsart}
\usepackage[margin=19mm]{geometry}
\usepackage{amsmath,amssymb,mathtools}
\usepackage[scr=rsfs,cal=euler]{mathalfa}
\usepackage{xfrac}
\usepackage[unicode,hidelinks,bookmarks=false]{hyperref}
\newcommand{\ie}{i.e.\ }
\newcommand{\cf}{cf.\ }
\newcommand{\resp}{respectively\ }
\newcommand{\Subsection}{\subsection}
\providecommand{\nobreakdash}{\nobreak\hbox{-}}
\setlength{\emergencystretch}{2em}
\multlinegap0pt
\usepackage{bm}
\usepackage{bbm}
\usepackage{dsfont}
\usepackage{xspace}
\usepackage{cancel}
\usepackage{mathtools}
\usepackage{amsthm}
\makeatletter
\@ifundefined{lemma}{\newtheorem{lemma}{Lemma}}{}
\makeatother
\newcommand{\gfnoa}[1]{F^{(#1)}(x)}
\title[Counting Colored Trees]{Counting Colored Trees}
\author{Stoyan Dimitrov, Nathan Fox, Kimberly Hadaway, Ashley Tharp, Stephan Wagner}
\date{Source: arXiv:2602.16055v2 (CC BY 4.0)}
\begin{document}
\maketitle

\noindent\emph{Source and licence.} Counting Colored Trees. Version arXiv:2602.16055v2 (CC BY 4.0), distributed under the Creative Commons Attribution 4.0 International licence, as recorded in the arXiv licence metadata for this exact version and shown on the abstract page https://arxiv.org/abs/2602.16055v2. Full rights notice: licenses/arxiv-2602.16055-CC-BY-4.0.txt.\par\smallskip

\noindent\textit{Editorial scope.} Counted unit: the Lemma whose source label is \texttt{lem:gen33} in the article (source0.tex lines 1711--1716, source label \texttt{lem:gen33}), delivered together with the article's own complete original proof of it (source0.tex lines 1717--1794, one \texttt{proof} environment of 78 lines). No mathematics is written, repaired or re-derived: statement and proof are the article's own text line for line. Required context retained uncounted: lem:gen33gf (lines 1683--1688). Every definition, display and label of those lines is retained unchanged. Omitted from this package and not redistributed: 1--1682, 1689--1710, 1795--5145. Nothing inside a retained excerpt is omitted or altered: every retained body is delivered exactly as the article's lines are written, so no omission receipt of any kind accompanies this package. Citations: the retained excerpts carry no citation command at all, so no bibliography entry of the article is transcribed and no reference is redistributed. Reference adaptation: none was needed and none was made; every reference inside the delivered unit resolves inside it, so no unresolved reference or citation is produced. Printed slips kept literal: no printed word, symbol, hypothesis or number of the counted statement or of its proof was altered. Layout and class: the article's own class is \texttt{amsart} with packages including a4wide, cite, tikz, hyperref, amsmath, amsthm, graphicx, standalone, mathdots, makecell, algorithm, array, ltablex, subcaption; the wrapper is the lane's pinned \texttt{amsart} editorial wrapper at 10pt on A4, which restates the theorem-like environments the retained text needs and adds the article's own macro definitions that the retained text needs (\texttt{\textbackslash gfnoa}), with extra wrapper packages (bm, bbm, dsfont, xspace, cancel, mathtools). The delivered numbering is the unit's own. Assets: no retained span contains a figure environment, a graphics-inclusion command, a PDF-page inclusion, a table or a TikZ picture; no image file is copied into this package, the package declares no asset dependency and no image file is redistributed with it. Licence: version arXiv:2602.16055v2 of this article is distributed under the Creative Commons Attribution 4.0 International licence, as recorded in the arXiv licence metadata for this exact version and shown on the abstract page https://arxiv.org/abs/2602.16055v2. A full rights notice is delivered at licenses/arxiv-2602.16055-CC-BY-4.0.txt. Characters: every retained excerpt is pure ASCII, so the delivered engine, XeLaTeX with its default Latin Modern text font, reports no missing character.
\begin{lemma}\label{lem:gen33gf}
    For all $m\geq1$, we have
    \[
    \gfnoa{m}\cdot\left(1-\sum_{i=1}^{m}\gfnoa{i}\right)=x.
    \]
\end{lemma}

\begin{lemma}\label{lem:gen33}
    For any $m\geq2$ and $1\leq j<m$, we have
    \[
    \gfnoa{m-j}=\frac{-\sqrt{x}P_{j+1}\!\left(\frac{\gfnoa{m}}{\sqrt{x}}\right)}{Q_{j+1}\!\left(\frac{\gfnoa{m}}{\sqrt{x}}\right)}+1-\sum_{i=1}^{m-j-1}\gfnoa{i}.
    \]
\end{lemma}

\begin{proof}
    We prove the lemma by induction on $j$. As a base case, we verify the statement when $j=1$. Starting from Lemma~\ref{lem:gen33gf}, we have
    \begin{align}
        &\nonumber \gfnoa{m}\cdot\left(1-\sum_{i=1}^{m}\gfnoa{i}\right)=x\\
        &\nonumber \Rightarrow\gfnoa{m}\cdot\left(1-\sum_{i=1}^{m-2}\gfnoa{i}-\gfnoa{m-1}-\gfnoa{m}\right)=x\\
        &\nonumber \Rightarrow1-\sum_{i=1}^{m-2}\gfnoa{i}-\gfnoa{m-1}-\gfnoa{m}=\frac{x}{\gfnoa{m}}\\
        &\nonumber \Rightarrow\gfnoa{m-1}=-\frac{x}{\gfnoa{m}}-\gfnoa{m}+1-\sum_{i=1}^{m-2}\gfnoa{i}\\
        &\Rightarrow\gfnoa{m-1}=\frac{-x-\gfnoa{m}^2}{\gfnoa{m}}+1-\sum_{i=1}^{m-2}\gfnoa{i}\label{eq:firstgfexpression}.
    \end{align}
    Next, we observe that $P_2(x)=P_1(x)^2+Q_1(x)^2=x^2+1$ and $Q_2(x)=P_1(x)\cdot Q_1(x)=x$. This implies
    \[
    P_2\!\left(\frac{\gfnoa{m}}{\sqrt{x}}\right)=\frac{\gfnoa{m}^2}{x}+1
    \]
    and
    \[
    Q_2\!\left(\frac{\gfnoa{m}}{\sqrt{x}}\right)=\frac{\gfnoa{m}}{\sqrt{x}}.
    \]
    By substitution followed by algebraic manipulation, we have 
    \begin{align}
    \frac{-\sqrt{x}P_{2}\!\left(\frac{\gfnoa{m}}{\sqrt{x}}\right)}{Q_{2}\!\left(\frac{\gfnoa{m}}{\sqrt{x}}\right)}&=\frac{-\sqrt{x}\cdot\left(\frac{\gfnoa{m}^2}{x}+1\right)}{\frac{\gfnoa{m}}{\sqrt{x}}}
    =\frac{-x\cdot\left(\frac{\gfnoa{m}^2}{x}+1\right)}{\gfnoa{m}}
    =\frac{-x-\gfnoa{m}^2}{\gfnoa{m}}\label{eq:secondgfexpression}.
    \end{align}
    Thus, we have, by substituting the last expression from~\eqref{eq:secondgfexpression} into~\eqref{eq:firstgfexpression}, 
    \[
    \gfnoa{m-1}=\frac{-\sqrt{x}P_{2}\!\left(\frac{\gfnoa{m}}{\sqrt{x}}\right)}{Q_{2}\!\left(\frac{\gfnoa{m}}{\sqrt{x}}\right)}+1-\sum_{i=1}^{m-2}\gfnoa{i},
    \]
    as required.

    For ease of notation in the inductive step, define
    \[
    R_j=\frac{-\sqrt{x}P_{j+1}\!\left(\frac{\gfnoa{m}}{\sqrt{x}}\right)}{Q_{j+1}\!\left(\frac{\gfnoa{m}}{\sqrt{x}}\right)}.
    \]
    For some $j\geq2$, suppose inductively that
    \[
    \gfnoa{m-j+1}=R_{j-1}+1-\sum_{i=1}^{m-j}\gfnoa{i}.
    \]
    By Lemma~\ref{lem:gen33gf}, we have
    \[
    \gfnoa{m-j+1}\cdot\left(1-\sum_{i=1}^{m-j+1}\gfnoa{i}\right)=x,
    \]
    which implies that
    \[
    \frac{x}{\gfnoa{m-j+1}}=1-\sum_{i=1}^{m-j+1}\gfnoa{i}=1-\sum_{i=1}^{m-j}\gfnoa{i}-\gfnoa{m-j+1}.
    \]
    By the inductive hypothesis, we have
    \[
    \frac{x}{R_{j-1}+1-\sum_{i=1}^{m-j}\gfnoa{i}}=1-\sum_{i=1}^{m-j}\gfnoa{i}-\left(R_{j-1}+1-\sum_{i=1}^{m-j}\gfnoa{i}\right)=-R_{j-1}.
    \]
    Rearranging terms, we see that
    \[
    -\frac{x}{R_{j-1}}=R_{j-1}+1-\sum_{i=1}^{m-j}\gfnoa{i}=R_{j-1}+1-\sum_{i=1}^{m-j-1}\gfnoa{i}-\gfnoa{m-j},
    \]
    and thus,
    \[
    \gfnoa{m-j}=R_{j-1}+\frac{x}{R_{j-1}}+1-\sum_{i=1}^{m-j-1}\gfnoa{i}.
    \]

    To conclude the proof, it suffices to show that
    \[
    R_j=R_{j-1}+\frac{x}{R_{j-1}}=\frac{R_{j-1}^2+x}{R_{j-1}}.
    \]
    We have
    \begin{align*}
    R_{j-1}^2+x&=\left(\frac{-\sqrt{x}P_j\!\left(\frac{\gfnoa{m}}{\sqrt{x}}\right)}{Q_{j}\!\left(\frac{\gfnoa{m}}{\sqrt{x}}\right)}\right)^2+x
    \\
    &=\frac{xP_{j}\!\left(\frac{\gfnoa{m}}{\sqrt{x}}\right)^2}{Q_{j}\!\left(\frac{\gfnoa{m}}{\sqrt{x}}\right)^2}+\frac{xQ_{j}\!\left(\frac{\gfnoa{m}}{\sqrt{x}}\right)^2}{Q_{j}\!\left(\frac{\gfnoa{m}}{\sqrt{x}}\right)^2}
    =\frac{xP_{j+1}\left(\frac{\gfnoa{m}}{\sqrt{x}}\right)}{Q_{j}\!\left(\frac{\gfnoa{m}}{\sqrt{x}}\right)^2}.
    \end{align*}
    Dividing through by $R_{j-1}$ gives
    \begin{align*}
        \frac{R_{j-1}^2+x}{R_{j-1}}&=\frac{xP_{j+1}\left(\frac{\gfnoa{m}}{\sqrt{x}}\right)}{Q_{j}\!\left(\frac{\gfnoa{m}}{\sqrt{x}}\right)^2}\cdot\frac{Q_{j}\!\left(\frac{\gfnoa{m}}{\sqrt{x}}\right)}{-\sqrt{x}P_j\!\left(\frac{\gfnoa{m}}{\sqrt{x}}\right)}
        =\frac{xP_{j+1}\left(\frac{\gfnoa{m}}{\sqrt{x}}\right)}{-\sqrt{x}P_{j}\!\left(\frac{\gfnoa{m}}{\sqrt{x}}\right)Q_{j}\!\left(\frac{\gfnoa{m}}{\sqrt{x}}\right)}\\
        &=\frac{-\sqrt{x}P_{j+1}\!\left(\frac{\gfnoa{m}}{\sqrt{x}}\right)}{Q_{j+1}\!\left(\frac{\gfnoa{m}}{\sqrt{x}}\right)}
        =R_j,
    \end{align*}
    as required.
\end{proof}

\end{document}
